\documentclass[10pt,a4paper]{amsart}
\usepackage{amsmath,amssymb,amsthm,graphicx}
% Retained preamble for the Stacks Project chapter {CH} (shared project preamble).
% The chapter file pinned-original-source.tex only carries \input{{preamble}} on its first line,
% so no preamble text lives inside the chapter itself.  The Stacks Project uses one single shared
% preamble for every chapter: the snapshot stages it three times (sources/stacks-chow/preamble.tex,
% sources/stacks-etale-cohomology/preamble.tex and focused-seeds/AGP/sources/stacks-snapshot-20261011/
% preamble.tex) and those copies are identical in size.  The block below is copied mechanically,
% line for line, from sources/stacks-chow/preamble.tex lines 14-18 and 154-217: the Xy-pic setup and
% the native theorem declarations together with the mathematical macros.  Nothing was added, removed
% or redefined; the wrapper class and the amsmath/amssymb/amsthm/graphicx base packages are supplied
% by the assembler.
% For commutative diagrams we use Xy-pic
% Chapter: stacks-spaces-more-morphisms.tex -- More on Morphisms of Spaces.
\usepackage[all]{xy}

% We use 2cell for 2-commutative diagrams.
\xyoption{2cell}

% Theorem environments.
%
\theoremstyle{plain}
\newtheorem{theorem}[subsection]{Theorem}
\newtheorem{proposition}[subsection]{Proposition}
\newtheorem{lemma}[subsection]{Lemma}

\theoremstyle{definition}
\newtheorem{definition}[subsection]{Definition}
\newtheorem{example}[subsection]{Example}
\newtheorem{exercise}[subsection]{Exercise}
\newtheorem{situation}[subsection]{Situation}

\theoremstyle{remark}
\newtheorem{remark}[subsection]{Remark}
\newtheorem{remarks}[subsection]{Remarks}

\numberwithin{equation}{subsection}

% Macros
%
\def\lim{\mathop{\mathrm{lim}}\nolimits}
\def\colim{\mathop{\mathrm{colim}}\nolimits}
\def\Spec{\mathop{\mathrm{Spec}}}
\def\Hom{\mathop{\mathrm{Hom}}\nolimits}
\def\Ext{\mathop{\mathrm{Ext}}\nolimits}
\def\SheafHom{\mathop{\mathcal{H}\!\mathit{om}}\nolimits}
\def\SheafExt{\mathop{\mathcal{E}\!\mathit{xt}}\nolimits}
\def\Sch{\mathit{Sch}}
\def\Mor{\mathop{\mathrm{Mor}}\nolimits}
\def\Ob{\mathop{\mathrm{Ob}}\nolimits}
\def\Sh{\mathop{\mathit{Sh}}\nolimits}
\def\NL{\mathop{N\!L}\nolimits}
\def\CH{\mathop{\mathrm{CH}}\nolimits}
\def\proetale{{pro\text{-}\acute{e}tale}}
\def\etale{{\acute{e}tale}}
\def\QCoh{\mathit{QCoh}}
\def\Ker{\mathop{\mathrm{Ker}}}
\def\Im{\mathop{\mathrm{Im}}}
\def\Coker{\mathop{\mathrm{Coker}}}
\def\Coim{\mathop{\mathrm{Coim}}}

% Boxtimes
%
\DeclareMathSymbol{\boxtimes}{\mathbin}{AMSa}{"02}

%
% Macros for moduli stacks/spaces
%
\def\QCohstack{\mathcal{QC}\!\mathit{oh}}
\def\Cohstack{\mathcal{C}\!\mathit{oh}}
\def\Spacesstack{\mathcal{S}\!\mathit{paces}}
\def\Quotfunctor{\mathrm{Quot}}
\def\Hilbfunctor{\mathrm{Hilb}}
\def\Curvesstack{\mathcal{C}\!\mathit{urves}}
\def\Polarizedstack{\mathcal{P}\!\mathit{olarized}}
\def\Complexesstack{\mathcal{C}\!\mathit{omplexes}}
% \Pic is the operator that assigns to X its picard group, usage \Pic(X)
% \Picardstack_{X/B} denotes the Picard stack of X over B
% \Picardfunctor_{X/B} denotes the Picard functor of X over B
\def\Pic{\mathop{\mathrm{Pic}}\nolimits}
\def\Picardstack{\mathcal{P}\!\mathit{ic}}
\def\Picardfunctor{\mathrm{Pic}}
\def\Deformationcategory{\mathcal{D}\!\mathit{ef}}
\usepackage{cleveref}
\crefname{originallemma}{original lemma}{original lemmas}
\Crefname{originallemma}{Original lemma}{Original lemmas}
\makeatletter
\newlabel{lemma-case-of-tor-independence}{{0CTQ}{source}{}{sourcepointer-1}{}}
\newlabel{lemma-case-of-tor-independence@cref}{{[originallemma][0][]0CTQ}{[1][0][]source}{}{}{}}
\makeatother
\crefname{originalreference}{original reference}{original references}
\Crefname{originalreference}{Original reference}{Original references}
\makeatletter
\newlabel{sites-cohomology-lemma-pseudo-coherent-pullback}{{08H4}{source}{}{sourcepointer-2}{}}
\newlabel{sites-cohomology-lemma-pseudo-coherent-pullback@cref}{{[originalreference][0][]08H4}{[1][0][]source}{}{}{}}
\makeatother
\crefname{originalreference}{original reference}{original references}
\Crefname{originalreference}{Original reference}{Original references}
\makeatletter
\newlabel{spaces-limits-lemma-descend-finite-presentation}{{07SK}{source}{}{sourcepointer-3}{}}
\newlabel{spaces-limits-lemma-descend-finite-presentation@cref}{{[originalreference][0][]07SK}{[1][0][]source}{}{}{}}
\makeatother
\crefname{originalreference}{original reference}{original references}
\Crefname{originalreference}{Original reference}{Original references}
\makeatletter
\newlabel{spaces-limits-lemma-descend-separated-morphism}{{0851}{source}{}{sourcepointer-4}{}}
\newlabel{spaces-limits-lemma-descend-separated-morphism@cref}{{[originalreference][0][]0851}{[1][0][]source}{}{}{}}
\makeatother
\crefname{originalreference}{original reference}{original references}
\Crefname{originalreference}{Original reference}{Original references}
\makeatletter
\newlabel{spaces-morphisms-lemma-scheme-theoretic-image-of-partial-section}{{0CP3}{source}{}{sourcepointer-5}{}}
\newlabel{spaces-morphisms-lemma-scheme-theoretic-image-of-partial-section@cref}{{[originalreference][0][]0CP3}{[1][0][]source}{}{}{}}
\makeatother
\crefname{originalreference}{original reference}{original references}
\Crefname{originalreference}{Original reference}{Original references}
\makeatletter
\newlabel{spaces-perfect-lemma-compare-base-change}{{08IR}{source}{}{sourcepointer-6}{}}
\newlabel{spaces-perfect-lemma-compare-base-change@cref}{{[originalreference][0][]08IR}{[1][0][]source}{}{}{}}
\makeatother
\crefname{originalreference}{original reference}{original references}
\Crefname{originalreference}{Original reference}{Original references}
\makeatletter
\newlabel{spaces-perfect-lemma-identify-pseudo-coherent-noetherian}{{08IK}{source}{}{sourcepointer-7}{}}
\newlabel{spaces-perfect-lemma-identify-pseudo-coherent-noetherian@cref}{{[originalreference][0][]08IK}{[1][0][]source}{}{}{}}
\makeatother
\title{Derived Chow's lemma for a separated algebraic space of finite presentation}
\author{The Stacks Project Authors}
\date{stacks-project:spaces-more-morphisms.tex:local-snapshot-20261011; GNU FDL 1.2 or later}
\begin{document}
\maketitle
\noindent\textit{Editorial note.} Counted span is the author's own proof environment, source lines 12033-12112, immediately following the statement at 12012-12031.\par\medskip
\noindent\textit{Editorial note.} Same-name pair note: this is the more general (algebraic spaces) side; the schemes side is not counted beside it.\par\medskip
\noindent\textit{Editorial note.} Six cross-chapter pointers (Sites and Sheaves, Limits of Spaces, Morphisms of Spaces, Derived Categories of Spaces) plus one same-chapter lemma are recorded in assembly.external\_references.\par\medskip
\noindent\textit{Allowed standard prerequisites (editorial).}\begin{itemize}
\item Descending finite presentation and separatedness along a limit. Limits of Spaces, lemmas spaces-limits-lemma-descend-finite-presentation and spaces-limits-lemma-descend-separated-morphism, cited in the proof. Moves the situation to a Noetherian approximation where the construction can be done.
\item Scheme-theoretic image of a partial section. Morphisms of Spaces, lemma spaces-morphisms-lemma-scheme-theoretic-image-of-partial-section, cited in the proof. Produces the closed subspace Z inside X x\_A P\textasciicircum{}n\_A with the required property over the open V.
\item Identifying a pushforward as pseudo-coherent in the Noetherian case. Derived Categories of Spaces, lemmas spaces-perfect-lemma-compare-base-change and spaces-perfect-lemma-identify-pseudo-coherent-noetherian, cited in the proof. Supplies the object E and its pseudo-coherence, which is what the derived statement adds to Chow's lemma.
\item Pseudo-coherence under pullback. Sites and Sheaves, lemma sites-cohomology-lemma-pseudo-coherent-pullback, cited in the proof. Keeps the pseudo-coherence conclusion stable along the base changes used in the construction.
\item Tor independence in the base change square. chapter lemma lemma-case-of-tor-independence, cited in the proof. Controls the derived tensor products appearing in the construction of b and c.
\end{itemize}\medskip
\ifdefined\hypertarget\hypertarget{sourcepointer-1}{}\fi
\noindent\textit{Source reference.} Original lemma 0CTQ: sources/stacks-spaces-more-morphisms/pinned-original-source.tex, source lines 11969--11969. Navigation pointer only: the counted proof cites this statement by \textbackslash{}ref as a standing input; the deduction the counted proof has to make is carried out inside the counted span, and no part of that deduction is delegated to this pointer.\par\medskip
\ifdefined\hypertarget\hypertarget{sourcepointer-2}{}\fi
\noindent\textit{Source reference.} Original reference 08H4: focused-seeds/geometry/S9-spaces-derived-chow-spaces/ref-tags.txt, source lines 5--5. Navigation pointer only: the counted proof cites this statement by \textbackslash{}ref as a standing input; the deduction the counted proof has to make is carried out inside the counted span, and no part of that deduction is delegated to this pointer.\par\medskip
\ifdefined\hypertarget\hypertarget{sourcepointer-3}{}\fi
\noindent\textit{Source reference.} Original reference 07SK: focused-seeds/geometry/S9-spaces-derived-chow-spaces/ref-tags.txt, source lines 6--6. Navigation pointer only: the counted proof cites this statement by \textbackslash{}ref as a standing input; the deduction the counted proof has to make is carried out inside the counted span, and no part of that deduction is delegated to this pointer.\par\medskip
\ifdefined\hypertarget\hypertarget{sourcepointer-4}{}\fi
\noindent\textit{Source reference.} Original reference 0851: focused-seeds/geometry/S9-spaces-derived-chow-spaces/ref-tags.txt, source lines 7--7. Navigation pointer only: the counted proof cites this statement by \textbackslash{}ref as a standing input; the deduction the counted proof has to make is carried out inside the counted span, and no part of that deduction is delegated to this pointer.\par\medskip
\ifdefined\hypertarget\hypertarget{sourcepointer-5}{}\fi
\noindent\textit{Source reference.} Original reference 0CP3: focused-seeds/geometry/S9-spaces-derived-chow-spaces/ref-tags.txt, source lines 8--8. Navigation pointer only: the counted proof cites this statement by \textbackslash{}ref as a standing input; the deduction the counted proof has to make is carried out inside the counted span, and no part of that deduction is delegated to this pointer.\par\medskip
\ifdefined\hypertarget\hypertarget{sourcepointer-6}{}\fi
\noindent\textit{Source reference.} Original reference 08IR: focused-seeds/geometry/S9-spaces-derived-chow-spaces/ref-tags.txt, source lines 9--9. Navigation pointer only: the counted proof cites this statement by \textbackslash{}ref as a standing input; the deduction the counted proof has to make is carried out inside the counted span, and no part of that deduction is delegated to this pointer.\par\medskip
\ifdefined\hypertarget\hypertarget{sourcepointer-7}{}\fi
\noindent\textit{Source reference.} Original reference 08IK: focused-seeds/geometry/S9-spaces-derived-chow-spaces/ref-tags.txt, source lines 10--10. Navigation pointer only: the counted proof cites this statement by \textbackslash{}ref as a standing input; the deduction the counted proof has to make is carried out inside the counted span, and no part of that deduction is delegated to this pointer.\par\medskip
\begin{lemma}[Derived Chow's lemma]
\label{lemma-derived-chow}
Let $A$ be a ring. Let $X$ be a separated algebraic space
of finite presentation over $A$. Let $x \in |X|$. Then there exist
an $n \geq 0$,
a closed subspace $Z \subset X \times_A \mathbf{P}^n_A$,
a point $z \in |Z|$,
an open $V \subset \mathbf{P}^n_A$, and
an object $E$ in $D(\mathcal{O}_{X \times_A \mathbf{P}^n_A})$ such that
\begin{enumerate}
\item $Z \to X \times_A \mathbf{P}^n_A$ is of finite presentation,
\item $c : Z \to \mathbf{P}^n_A$ is a closed immersion over $V$,
set $W = c^{-1}(V)$,
\item the restriction of $b : Z \to X$ to $W$ is \'etale,
$z \in W$, and $b(z) = x$,
\item $E|_{X \times_A V} \cong
(b, c)_*\mathcal{O}_Z|_{X \times_A V}$,
\item $E$ is pseudo-coherent and supported on $Z$.
\end{enumerate}
\end{lemma}

\begin{proof}
We can find a finite type $\mathbf{Z}$-subalgebra $A' \subset A$
and an algebraic space $X'$ separated and of finite presentation over $A'$
whose base change to $A$ is $X$. See
Limits of Spaces, Lemmas
\ref{spaces-limits-lemma-descend-finite-presentation} and
\ref{spaces-limits-lemma-descend-separated-morphism}.
Let $x' \in |X'|$ be the image of $x$.
If we can prove the lemma for $(X'/A', x')$, then
the lemma follows for $(X/A, x)$.
Namely, if $n', Z', z', V', E'$ provide the solution
for $(X'/A', x')$, then we can let
$n = n'$,
let $Z \subset X \times \mathbf{P}^n$ be the inverse image of $Z'$,
let $z \in Z$ be the unique point mapping to $x$,
let $V \subset \mathbf{P}^n_A$ be the inverse image of $V'$, and
let $E$ be the derived pullback of $E'$.
Observe that $E$ is pseudo-coherent by
Cohomology on Sites, Lemma
\ref{sites-cohomology-lemma-pseudo-coherent-pullback}.
It only remains to check (5). To see this
set $W = c^{-1}(V)$ and $W' = (c')^{-1}(V')$
and consider the cartesian square
$$
\xymatrix{
W \ar[d]_{(b, c)} \ar[r] & W' \ar[d]^{(b', c')} \\
X \times_A V \ar[r] & X' \times_{A'} V'
}
$$
By Lemma \ref{lemma-case-of-tor-independence} $X \times_A V$ and $W'$
are tor-independent over $X' \times_{A'} V'$.
Thus the derived pullback of
$(b', c')_*\mathcal{O}_{W'}$ to $X \times_A V$
is $(b, c)_*\mathcal{O}_W$ by
Derived Categories of Spaces,
Lemma \ref{spaces-perfect-lemma-compare-base-change}.
This also uses that $R(b', c')_*\mathcal{O}_{Z'} = (b', c')_*\mathcal{O}_{Z'}$
because $(b', c')$ is a closed immersion and similarly for
$(b, c)_*\mathcal{O}_Z$.
Since $E'|_{U' \times_{A'} V'} =
(b', c')_*\mathcal{O}_{W'}$ we obtain
$E|_{U \times_A V} = (b, c)_*\mathcal{O}_W$
and (5) holds.
This reduces us to the situation described in the next
paragraph.

\medskip\noindent
Assume $A$ is of finite type over $\mathbf{Z}$.
Choose an \'etale morphism $U \to X$ where $U$ is an affine scheme
and a point $u \in U$ mapping to $x$. Then $U$ is of finite type over $A$.
Choose a closed immersion $U \to \mathbf{A}^n_A$ and denote
$j : U \to \mathbf{P}^n_A$ the immersion we get by composing
with the open immersion $\mathbf{A}^n_A \to \mathbf{P}^n_A$.
Let $Z$ be the scheme theoretic closure of
$$
(\text{id}_U, j) : U \longrightarrow X \times_A \mathbf{P}^n_A
$$
Let $z \in Z$ be the image of $u$.
Let $Y \subset \mathbf{P}^n_A$ be the scheme theoretic
closure of $j$. Then it is clear that $Z \subset X \times_A Y$
is the scheme theoretic closure of
$(\text{id}_U, j) : U \to X \times_A Y$.
As $X$ is separated, the morphism
$X \times_A Y \to Y$ is separated as well.
Hence we see that $Z \to Y$ is an isomorphism over
the open subscheme $j(U) \subset Y$ by
Morphisms of Spaces, Lemma
\ref{spaces-morphisms-lemma-scheme-theoretic-image-of-partial-section}.
Choose $V \subset \mathbf{P}^n_A$ open with $V \cap Y = j(U)$.
Then we see that (2) holds, that $W = (\text{id}_U, j)(U)$, and hence
that (3) holds. Part (1) holds because $A$ is Noetherian.

\medskip\noindent
Because $A$ is Noetherian we see that $X$ and $X \times_A \mathbf{P}^n_A$
are Noetherian algebraic spaces. Hence we can take $E = (b, c)_*\mathcal{O}_Z$
in this case: (4) is clear and for (5) see
Derived Categories of Spaces, Lemma
\ref{spaces-perfect-lemma-identify-pseudo-coherent-noetherian}.
This finishes the proof.
\end{proof}

\end{document}
